\documentclass[11pt,a4paper,sans]{moderncv}

% --- стиль и цвет ---
\moderncvstyle{classic}
\moderncvcolor{blue}

% --- шрифты и язык (по заданию) ---
\usepackage{fontspec}
\defaultfontfeatures{Renderer=Basic, Ligatures=TeX}
\newfontfamily\cyrillicfonttt{CMU Typewriter Text}
\newfontfamily\cyrillicfont{CMU Sans Serif}
\setmainfont{CMU Sans Serif}
\setsansfont{CMU Sans Serif}
\setmonofont{CMU Typewriter Text}

\usepackage{polyglossia}
\setdefaultlanguage{russian}
\setotherlanguages{english}

% --- поля страницы ---
\usepackage[scale=0.8]{geometry}

% --- личные данные ---
\name{Андрей}{Фролов}
\title{Студент 3 курса, ИВТ}
\address{Санкт-Петербург, Россия}{Родной город~--- Ярославль}{}
\phone[mobile]{+7~(000)~000~00~00}
\email{andrewkafrolov76@gmail.com}
\social[github]{bluv4nul}
\social[telegram]{bluv4nul}
\photo[64pt][0.4pt]{avatar}

\begin{document}
\makecvtitle

\section{О себе}
\cvitem{}{Интересуюсь бэкенд-разработкой и работой с данными. Пишу REST API
  на Python (FastAPI) и Node.js (Express), работаю с реляционными и документными
  базами данных. Ищу стажировку или практику на позиции backend-разработчика
  с возможностью совмещать с учёбой.}

\section{Образование}
\cventry{2024--2028}
  {Бакалавриат, 09.03.01 «Информатика и вычислительная техника»}
  {Российский государственный педагогический университет им.~А.\,И.~Герцена}
  {Санкт-Петербург}
  {}
  {Профиль «Технологии разработки программного обеспечения и обработки больших данных».
   Институт информационных технологий и технологического образования.}

\section{Навыки}
\cvitem{Языки}{Python, JavaScript (Node.js), SQL}
\cvitem{Бэкенд}{FastAPI, Pydantic, SQLAlchemy, Express}
\cvitem{Фронтенд}{Vue~3, Vite, Vue Router, Axios}
\cvitem{Базы данных}{MySQL, SQLite, MongoDB}
\cvitem{Инструменты}{Git, Docker, Docker Compose}
\cvitem{Облака}{Yandex Cloud (ознакомительно)}

\section{Иностранные языки}
\cvitemwithcomment{Русский}{родной}{}
\cvitemwithcomment{Английский}{B2}{}

\section{Проекты}
\cventry{}
  {ElectroLibrary~--- электронная библиотека (командный проект)}
  {FastAPI, MongoDB, Vue~3, Docker Compose}
  {\href{https://github.com/bluv4nul/Study-Library-Project}{github.com/bluv4nul/Study-Library-Project}}
  {}
  {Fullstack-приложение из трёх сервисов: REST API на FastAPI, SPA на Vue~3,
   MongoDB; запуск через Docker Compose.
   \textbf{Моя роль:} база данных~--- структура коллекций MongoDB,
   init-скрипты начального наполнения, подключение к API через PyMongo.}

\cventry{}
  {REST API курсов валют ЦБ РФ}
  {FastAPI, SQLAlchemy, Pydantic, SQLite}
  {\href{https://github.com/bluv4nul/cbrrf-rest-api}{github.com/bluv4nul/cbrrf-rest-api}}
  {}
  {Асинхронный API: получение курсов с сайта ЦБ РФ, CRUD пользователей,
   подписки на валюты (связи «один ко многим»), валидация через Pydantic.
   Сборка deb-пакета.}

\cventry{}
  {REST API «Калькулятор»}
  {Node.js, Express}
  {\href{https://github.com/bluv4nul/frolov-calc}{github.com/bluv4nul/frolov-calc}}
  {}
  {Арифметические эндпоинты, валидация параметров через middleware,
   ответы в JSON, обработка ошибок (код 400, деление на ноль).
   Развёрнут на \href{https://frolov-calc.ru}{frolov-calc.ru}.}

\end{document}